% =====================================================================
% MARCO ETICO Y LEGAL -- Responsable: Fabian Moreno Ugarte
%
% REGLA DE REDACCION APLICADA EN ESTE ARCHIVO:
%   - Ninguna cita normativa se escribe de memoria. Donde el numeral, la
%     fecha o la redaccion literal no se pudieron contrastar contra la
%     fuente primaria, el texto lo dice expresamente y la afirmacion se
%     formula en condicional o se atribuye al encargo del proyecto. Los
%     puntos de verificacion abiertos estan en PENDIENTES.md.
%   - No se afirman conclusiones juridicas categoricas. Las zonas de
%     interpretacion se redactan como riesgos a validar.
%   - La trazabilidad de cada afirmacion esta en NOTAS_FUENTES.md.
% =====================================================================
\section{Marco ético y legal aplicable}
\label{sec:marco}

\subsection{Objeto y método de esta sección}
\label{sec:marco-metodo}

Esta sección identifica el marco normativo peruano aplicable al
tratamiento de imágenes de personas mediante videovigilancia y visión
computacional en un terminal aeroportuario, y evalúa en qué medida la
solución propuesta por el equipo se ajusta a él.

El método seguido tiene tres pasos. Primero se establece la jerarquía
normativa aplicable, de la Constitución hacia abajo. Segundo, se traduce
cada obligación en un requisito verificable sobre el sistema. Tercero, se
señalan las zonas en las que la norma admite más de una lectura
razonable, que se presentan como riesgos a validar y no como conclusiones.

\cajaaviso{Advertencia metodológica}{Este análisis lo elabora un equipo
de estudiantes de ingeniería, no un despacho legal. Su valor es
identificar y ordenar las obligaciones aplicables y traducirlas a
requisitos técnicos. La calificación jurídica definitiva ---
particularmente la de si un dato concreto constituye o no dato sensible,
y la de qué base de legitimación ampara cada tratamiento --- corresponde
al área legal de LAP y, en su caso, a la Autoridad Nacional de Protección
de Datos Personales (ANPD). Donde este documento afirma que \enquote{es
razonable sostener} algo, debe leerse exactamente así y no como un
dictamen.}

\subsection{Nivel constitucional}
\label{sec:constitucional}

El artículo 2 de la Constitución Política del Perú de 1993 reconoce, como
derechos fundamentales de la persona, dos incisos directamente
pertinentes al proyecto.

\begin{description}[leftmargin=0pt,labelindent=0pt,style=unboxed,itemsep=0.5em]
  \item[Inciso 6 --- autodeterminación informativa.] Reconoce el derecho
    a que los servicios informáticos, computarizados o no, públicos o
    privados, no suministren informaciones que afecten la intimidad
    personal y familiar. Este inciso es la base constitucional del
    régimen de protección de datos personales y el fundamento del proceso
    de \emph{hábeas data}. Su relevancia para el proyecto es que sitúa el
    control sobre la información personal como un derecho fundamental, no
    como una obligación meramente administrativa: un incumplimiento en
    esta materia no se agota en una multa.
  \item[Inciso 7 --- honor, intimidad, voz e imagen propias.] Reconoce el
    derecho al honor y a la buena reputación, a la intimidad personal y
    familiar, así como a la voz y a la imagen propias. Es el inciso más
    directamente aplicable a un sistema de videovigilancia, porque
    protege la imagen de la persona como tal.
\end{description}

Ambos incisos se recogen aquí en paráfrasis y no en su redacción literal.
La distinción importa para el uso que pueda darse a este texto: sirve
para situar el marco de derechos aplicable, pero no para citarse como
transcripción. La cita textual debe tomarse de la edición oficial de la
Constitución.

De estos dos incisos se desprende el criterio rector que atraviesa toda
esta sección: la captación de la imagen de una persona en un espacio
abierto al público no es intrínsecamente ilícita, pero constituye una
injerencia en un derecho fundamental y, por lo tanto, debe superar un
examen de proporcionalidad. La pregunta pertinente no es
\enquote{¿está prohibido?} sino \enquote{¿es idóneo, necesario y
proporcional para la finalidad declarada?}.

\subsection{Ley 29733 y su nuevo reglamento (D.S. 016-2024-JUS)}
\label{sec:ley29733}

\subsubsection{Estructura general}

La Ley 29733, Ley de Protección de Datos Personales, es la norma central
del régimen. Su reglamento vigente es el aprobado por Decreto Supremo
016-2024-JUS.

\begin{center}
\small
\begin{tabularx}{\linewidth}{@{}L{4.6cm}Y@{}}
\toprule
\textbf{Dato} & \textbf{Valor (según el encargo del proyecto)} \\
\midrule
Norma reglamentaria vigente & D.S. 016-2024-JUS \\
Fecha de publicación & 30 de noviembre de 2024 \\
Fecha de entrada en vigencia & 30 de marzo de 2025 \\
Norma derogada & D.S. 003-2013-JUS \\
\bottomrule
\end{tabularx}
\end{center}

El encabezado de la segunda columna debe leerse literalmente: los cuatro
datos provienen del encargo interno del proyecto y no han sido
contrastados contra el diario oficial \emph{El Peruano}. Se reproducen
como referencia de trabajo y no como dato verificado, y el cuadro no
debería salir del equipo sin esa validación previa, porque un error en la
fecha de entrada en vigencia arrastraría consigo todo el análisis de
régimen transitorio.

Si las fechas se confirman, la consecuencia práctica de la sucesión
normativa es que cualquier documentación de cumplimiento que LAP tenga
elaborada bajo el reglamento anterior debe revisarse. Los procedimientos, cláusulas informativas y
registros construidos bajo el D.S. 003-2013-JUS no son automáticamente
conformes con el nuevo reglamento.

\subsubsection{Datos biométricos como dato sensible}
\label{sec:biometricos}

La Ley 29733 define los datos sensibles como una categoría reforzada, e
incluye en ella los datos biométricos que por sí mismos puedan
identificar al titular, junto con los datos sobre origen racial y étnico,
ingresos económicos, convicciones políticas, religiosas, filosóficas o
morales, afiliación sindical, y la información relativa a la salud o a la
vida sexual.

Esta descripción se ofrece sin cita de numeral. El artículo 2 de la Ley
29733 es el que contiene la definición, pero el inciso exacto y la
eventual precisión que el D.S. 016-2024-JUS haya introducido sobre ella
no se han contrastado contra el texto vigente, de modo que el párrafo
anterior describe la categoría y no la transcribe.

Esta calificación importa porque el tratamiento de datos sensibles está
sujeto a un régimen más estricto que el de los datos personales comunes.
De ahí se sigue la pregunta técnica central de todo este informe:

\cajaaviso{La pregunta que decide el régimen aplicable}{¿El sistema
propuesto trata datos biométricos? La respuesta no depende de que haya
cámaras, sino de si en algún punto del \emph{pipeline} se extrae de la
imagen una característica que por sí sola permita identificar a una
persona --- una plantilla facial, un vector de embedding facial, un
patrón de marcha. Un mapa de densidad y un conjunto de coordenadas de
cabezas no cumplen esa condición. Sin embargo, el fotograma original del
que se derivan sí contiene rostros, y su tratamiento --- aunque sea
transitorio y en memoria --- forma parte del ciclo. La calificación
depende, por tanto, de la arquitectura concreta y no del modelo aislado.
La auditoría de la Sección~\ref{sec:auditoria} desarrolla este punto.}

Conviene ser explícito sobre un matiz frecuentemente pasado por alto: que
la salida del modelo no sea biométrica no significa que la entrada no lo
sea. Un fotograma de un hall de aeropuerto contiene rostros identificables
en calidad suficiente para reconocimiento facial. Si ese fotograma se
almacena, se transmite fuera del perímetro o se conserva para
reentrenamiento, el análisis de riesgo cambia sustancialmente aunque el
modelo desplegado nunca haga reconocimiento facial.

\subsubsection{Evaluación de Impacto en Protección de Datos (EIPD)}
\label{sec:eipd}

El D.S. 016-2024-JUS exige realizar una Evaluación de Impacto en
Protección de Datos \textbf{antes de iniciar} un tratamiento de alto
riesgo. Entre los supuestos que activan la obligación figuran
expresamente la videovigilancia masiva, el perfilado automatizado
mediante inteligencia artificial y el tratamiento de datos biométricos a
escala. La evaluación debe documentarse y, según el nivel de riesgo
residual, puede requerir consulta previa a la Autoridad Nacional de
Protección de Datos Personales.

Conviene ser directo sobre lo que esto significa para el proyecto. El
sistema parte de la captación masiva y continua de imágenes de personas
en las zonas comunes de un terminal aeroportuario. Ese solo hecho ---
con independencia de que el modelo produzca conteos agregados y no
identifique a nadie --- sitúa al tratamiento dentro del supuesto de
videovigilancia masiva. No es una zona gris: es el supuesto típico que
la norma describe.

\cajaaviso{La EIPD es un entregable del proyecto, no un obstáculo a
esquivar}{Hay una tentación previsible en este punto: argumentar que,
como la salida es agregada y no hay reconocimiento facial, el sistema no
alcanza el umbral de alto riesgo y la EIPD no sería exigible. Este frente
recomienda \textbf{no} construir el argumento en esa dirección, por dos
razones.

La primera es que el supuesto que se activa es el de captación masiva, y
esa captación existe con independencia de lo que el modelo haga después.
La segunda es que la EIPD es precisamente el documento en el que las
decisiones de diseño de este informe --- minimización, no persistencia,
salida agregada, procesamiento local --- se registran como medidas de
mitigación y quedan acreditadas. Renunciar a la EIPD sería renunciar al
lugar donde el trabajo de este frente se vuelve exigible y demostrable
ante la autoridad.

En consecuencia, se propone tratar la EIPD como un entregable del
proyecto: no como un trámite del que el sistema se escapa, sino como el
producto que ordena y da respaldo formal a todo lo demás.}

La estructura de la Sección~\ref{sec:analisis} de este informe ---
identificación del tratamiento, valoración de necesidad y
proporcionalidad, identificación de riesgos para los derechos de las
personas y propuesta de medidas de mitigación --- corresponde
deliberadamente a la de una evaluación de impacto, y constituye el
insumo principal para redactarla.

Tres extremos de esta obligación quedan descritos pero no citados: el
artículo concreto del D.S. 016-2024-JUS que la regula, el contenido
mínimo que la norma exige al documento, y los supuestos en que la
consulta previa a la ANPD pasa de facultativa a obligatoria. La
existencia de la obligación está verificada; su articulado, no. Al
redactar la EIPD habrá que cerrar esos tres puntos, porque el contenido
mínimo condiciona la estructura del documento.

Queda abierta, además, una cuestión de alcance que el equipo debe zanjar
antes de comprometerse ante el cliente: si la EIPD se ofrece como
entregable formal del Capstone o si el proyecto aporta su estructura y su
contenido sustantivo dejando la formalización a LAP. La segunda opción es
más realista para el alcance del curso y la primera es más valiosa para
el cliente; lo que no admite ambigüedad es que se diga cuál de las dos se
está ofreciendo.

\subsubsection{Notificación de brechas de seguridad}
\label{sec:brechas}

Según el encargo del proyecto, el D.S. 016-2024-JUS introduce la
obligación de notificar las brechas de seguridad de datos personales
dentro de un plazo de 48 horas.

Ese plazo se toma del encargo y no del texto de la norma, y los detalles
que determinan cómo se cumple siguen sin contrastar: el artículo que lo
fija, el hecho desde el que se computa --- conocimiento de la brecha,
confirmación u otro ---, si son horas hábiles o calendario, si la
notificación se dirige a la ANPD, a los titulares afectados o a ambos, y
si existe un umbral de gravedad por debajo del cual no se notifica.
Ninguno de esos extremos debería darse por sabido al redactar el
procedimiento de respuesta.

Con independencia de esos detalles, la obligación tiene una consecuencia
de ingeniería que no puede resolverse
después del incidente: un plazo de 48 horas solo es cumplible si existe
de antemano capacidad de detectar la brecha, de determinar su alcance y
de identificar a los titulares afectados. Esto exige registro de
auditoría (\emph{logging}) sobre los accesos, inventario actualizado de
dónde reside cada categoría de dato, y un procedimiento de respuesta
ensayado. Un sistema sin trazas de acceso no puede, materialmente,
notificar en plazo lo que no puede reconstruir.

\begin{description}[leftmargin=0pt,style=unboxed,itemsep=0.35em]
  \item[Requisito técnico derivado R-01.] Registro de auditoría inmutable
    de todo acceso a fotogramas y a salidas del modelo, con
    identificación de usuario, marca temporal y recurso accedido.
  \item[Requisito técnico derivado R-02.] Inventario vigente de flujos y
    de ubicaciones de almacenamiento, actualizado ante cada cambio de
    arquitectura.
  \item[Requisito técnico derivado R-03.] Procedimiento documentado de
    respuesta a incidentes con responsable designado y ensayo periódico.
\end{description}

Este frente propone documentar R-01, R-02 y R-03 como \emph{requisitos de
puesta en producción} y no como entregables del ciclo. El sistema del
Capstone es un prototipo académico, y comprometer un procedimiento
ensayado de respuesta a incidentes excedería su alcance; en cambio,
dejarlos enunciados como condición previa al despliegue real sitúa la
obligación donde corresponde sin prometer lo que el equipo no puede
sostener.

\subsubsection{Oficial de Datos Personales}
\label{sec:odp}

El D.S. 016-2024-JUS incorpora la figura del Oficial de Datos Personales
(ODP), función equivalente en su lógica al \emph{Data Protection Officer}
del régimen europeo.

La obligación de designarlo no es uniforme ni inmediata: el reglamento
establece un \textbf{cronograma escalonado} en función de los ingresos
anuales de la empresa, de modo que las organizaciones de mayor facturación
quedan obligadas antes que las de menor tamaño.

Esto cambia la pregunta que corresponde llevar al cliente. No es si LAP
tiene o no un ODP, sino \textbf{en qué tramo del cronograma se ubica} y,
por tanto, desde qué fecha le resulta exigible la designación. Dado el
volumen de operación del concesionario del principal terminal aéreo del
país, lo esperable es que se encuentre en uno de los tramos iniciales,
aunque esto debe confirmarse y no asumirse.

La existencia del cronograma escalonado está verificada; sus términos
concretos, no. Quedan por contrastar el artículo que lo fija, los
umbrales de ingresos anuales que delimitan cada tramo con sus respectivas
fechas, la eventual obligación de comunicar la designación a la ANPD y
los requisitos de independencia o de ausencia de conflicto de interés que
se exijan al ODP. Por eso el párrafo anterior habla de lo esperable y no
de lo exigible.

Para el proyecto, la relevancia es práctica: cuando la designación sea
exigible, el ODP es el interlocutor natural del equipo. Todas las
decisiones de diseño con impacto en privacidad que este informe deja
abiertas deberían validarse con esa persona antes de un despliegue real,
y la EIPD de la Sección~\ref{sec:eipd} es un documento que ordinariamente
pasa por su revisión.

De ahí que este frente proponga incorporar dos preguntas al primer
contacto formal con el cliente: en qué tramo del cronograma se ubica LAP
y si ya designó al ODP. Si la designación aún no le resulta exigible, el
entregable de cumplimiento tendrá que dirigirse entretanto al área legal,
al área de seguridad de la información o a ambas, y conviene que ese
destinatario esté acordado antes de redactarlo.

\subsubsection{Alcance extraterritorial}
\label{sec:extraterritorial}

El D.S. 016-2024-JUS contempla, según el encargo del proyecto, un alcance
extraterritorial del régimen peruano de protección de datos.

El dato procede del encargo y no del texto de la norma: ni el artículo
que fija ese alcance ni sus criterios de conexión --- establecimiento en
el país, uso de medios situados en el Perú, oferta de bienes o servicios
a titulares en el Perú, monitoreo de comportamiento u otros --- han sido
contrastados. Lo que sigue vale, por tanto, como identificación de los
flujos que habría que analizar si el alcance se confirma, y no como
calificación jurídica de cada uno.

Este punto tiene una consecuencia directa y concreta sobre decisiones que
el equipo técnico probablemente ya tomó o está por tomar:

\begin{itemize}[leftmargin=1.5em,itemsep=0.3em]
  \item Si el entrenamiento o la inferencia se ejecutan en infraestructura
        de nube alojada fuera del Perú, hay un flujo transfronterizo que
        debe analizarse.
  \item Si se emplean servicios gestionados de terceros (etiquetado,
        \emph{MLOps}, almacenamiento de objetos), esos terceros son
        encargados de tratamiento y requieren contrato con las cláusulas
        que la norma exija.
  \item El uso de un modelo preentrenado descargado del exterior no
        constituye por sí mismo transferencia de datos personales, porque
        no se exportan datos; pero enviar fotogramas del terminal a una
        API externa para inferencia sí lo constituiría.
\end{itemize}

\cajaaviso{Riesgo a validar}{La distinción del último punto es
importante y suele confundirse. Descargar pesos de un modelo hacia el
Perú no exporta datos personales. Enviar imágenes del terminal hacia un
servicio en el extranjero sí. Si la arquitectura del equipo contempla
inferencia en la nube, el análisis de transferencia internacional deja de
ser hipotético. Esta es, a juicio de este frente, la decisión de
arquitectura con mayor impacto legal de todo el proyecto ---
\emph{on-premise} en el terminal o en nube externa ---, y al cerrar esta
versión del informe no consta por escrito. Debe confirmarse con el frente
técnico y documentarse.}

\subsection{Directiva 01-2020-JUS/DGTAIPD sobre videovigilancia}
\label{sec:directiva}

La Dirección General de Transparencia, Acceso a la Información Pública y
Protección de Datos Personales (DGTAIPD) emitió la Directiva
01-2020-JUS/DGTAIPD, referida al tratamiento de datos personales mediante
sistemas de videovigilancia.

La Directiva se encuentra vigente y resulta aplicable al proyecto. La
entrada en vigencia del D.S. 016-2024-JUS no la desplazó: el Congreso de
la República la invoca como base legal de su propio sistema de
videovigilancia en la Resolución 088-2025, posterior al nuevo
reglamento.\footnote{La confirmación definitiva exige revisar las
disposiciones derogatorias y complementarias del D.S. 016-2024-JUS, para
descartar una derogación parcial que no se refleje en el uso que otras
entidades hacen de la Directiva. Que un organismo del Estado la cite como
base legal vigente es un indicio sólido de su vigencia, pero no equivale
a la lectura del texto derogatorio.}

Lo que resta verificar de esta norma no es su vigencia sino su
identificación formal: la denominación oficial completa y el número y la
fecha de la resolución directoral que la aprueba. Hasta cerrar ese punto,
la Directiva se invoca aquí por su número y no por su título oficial.

Dos obligaciones de la Directiva estructuran el diseño del sistema.

\subsubsection{Deber de información mediante carteles}

El tratamiento de datos mediante videovigilancia exige informar
previamente a los titulares mediante carteles informativos ubicados de
forma visible en las zonas videovigiladas. La lógica es que el
consentimiento no es la base habilitante típica en videovigilancia de
seguridad, pero el deber de información subsiste con independencia de la
base que se invoque: la persona debe poder saber que está siendo captada,
por quién y para qué.

En el Anexo~\ref{anx:cartel} se propone un modelo de cartel para el
proyecto, construido sobre los elementos que la Directiva exige.

\subsubsection{Proporcionalidad}

La Directiva recoge el principio de proporcionalidad: la videovigilancia
debe ser idónea para la finalidad perseguida, necesaria en el sentido de
que no exista un medio menos invasivo igualmente eficaz, y proporcional
en sentido estricto.

Este principio es el que ofrece el argumento más sólido a favor del
enfoque del proyecto, y conviene formularlo con precisión porque será el
eje de la defensa del sistema ante el cliente:

\begin{center}
\small
\begin{tabularx}{\linewidth}{@{}L{3.0cm}Y@{}}
\toprule
\textbf{Subprincipio} & \textbf{Argumento aplicado al proyecto} \\
\midrule
Idoneidad &
La estimación automática de densidad permite detectar aglomeraciones que
la observación humana de decenas de monitores no detecta de forma
oportuna. El medio es apto para el fin. \\
\addlinespace
Necesidad &
Es el subprincipio decisivo. Si la finalidad es estimar aglomeraciones,
el medio menos invasivo que la alcanza es un sistema que produce conteos
agregados y no identifica personas. Un sistema con reconocimiento facial
alcanzaría la misma finalidad siendo más invasivo, y por eso no
superaría este examen. \\
\addlinespace
Proporcionalidad en sentido estricto &
El beneficio esperado --- prevención de incidentes de seguridad por
aglomeración y mejora del flujo de pasajeros --- debe ponderarse contra
la injerencia en la imagen de miles de personas al día. La ponderación
depende de la magnitud real del problema en el terminal. \\
\bottomrule
\end{tabularx}
\end{center}

Conviene señalar la debilidad de este razonamiento en lugar de
disimularla. El tercer subprincipio exige ponderar el beneficio esperado
contra la injerencia, y esa ponderación requiere evidencia sobre la
magnitud real del problema: incidentes por aglomeración registrados,
tiempos de cola observados, quejas recibidas. Solo LAP puede aportar esos
datos. Mientras no los aporte, el examen de proporcionalidad en sentido
estricto se sostiene sobre una premisa no acreditada, y así debe
presentarse ante el cliente y no como un juicio ya emitido.

\subsection{D.L. 1218 y Ley 30120}
\label{sec:dl1218}

Dos normas complementan el régimen en materia de videovigilancia asociada
a seguridad ciudadana:

\begin{description}[leftmargin=0pt,style=unboxed,itemsep=0.4em]
  \item[Decreto Legislativo 1218.] Regula el uso de las cámaras de
    videovigilancia. No se citan aquí ni su denominación oficial completa
    ni su fecha, y tampoco sus disposiciones sobre la obligación de
    entregar imágenes a la Policía Nacional del Perú y al Ministerio
    Público ni sobre los plazos de conservación que establezca: son
    justamente los extremos que condicionan la política de retención del
    sistema, y ninguno ha sido contrastado contra el texto de la norma.
  \item[Ley 30120.] Ley de apoyo a la seguridad ciudadana con cámaras de
    videovigilancia en áreas de dominio público. Queda por determinar su
    ámbito de aplicación respecto de establecimientos privados abiertos
    al público, que es lo que decide si alcanza o no al terminal. Este
    informe no lo afirma en ningún sentido.
\end{description}

\cajaaviso{Zona gris identificada --- naturaleza del espacio}{Un terminal
aeroportuario concesionado a un operador privado no encaja limpiamente ni
en la categoría de área de dominio público ni en la de recinto privado
cerrado. Es un espacio de titularidad y gestión privada, abierto al
público, con funciones de interés público y sujeto a regulación sectorial
aeronáutica.

De esa calificación depende qué régimen de videovigilancia resulta
exigible y con qué intensidad. Este informe \textbf{no resuelve} la
cuestión, porque hacerlo requeriría un análisis del contrato de concesión
y de la normativa sectorial que excede tanto el alcance del Capstone como
la competencia del equipo. Se deja planteada como el primer punto a
elevar al área legal de LAP, junto con la determinación del régimen
jurídico concreto del Aeropuerto Internacional Jorge Chávez y de la
normativa aeronáutica que le resulte aplicable en materia de vigilancia.}

Existe además una tensión que conviene anticipar. Estas normas de
seguridad ciudadana pueden generar obligaciones de conservación y de
entrega de imágenes a autoridades, mientras que el principio de
minimización de la Ley 29733 empuja hacia conservar lo mínimo por el
menor tiempo posible. No es una contradicción insalvable --- las
finalidades son distintas y pueden coexistir con bases y plazos separados
--- pero sí obliga a que la política de retención distinga explícitamente
entre el video de seguridad y los datos derivados del sistema de conteo.

\cajaaviso{Recomendación de diseño}{Tratar el sistema de estimación de
aglomeraciones como un subsistema \textbf{separado} del CCTV de
seguridad, con su propia finalidad, su propia base de legitimación y su
propio plazo de retención. Mezclarlos haría que el subsistema de conteo
herede las obligaciones de conservación del CCTV de seguridad, lo que
perjudica el argumento de minimización sin aportar beneficio operativo.}

\subsection{Referencia comparada}
\label{sec:comparada}

Las normas que siguen \textbf{no son aplicables} al proyecto por
territorio. Se incorporan porque constituyen el estándar de referencia
internacional en la materia y porque adoptarlas voluntariamente refuerza
la posición de LAP; en ningún caso sustituyen al régimen peruano.

\subsubsection{RGPD, artículo 35 --- evaluación de impacto}

El artículo 35 del Reglamento (UE) 2016/679 exige una evaluación de
impacto relativa a la protección de datos cuando un tratamiento entrañe
un alto riesgo para los derechos y libertades de las personas físicas,
y menciona expresamente entre esos supuestos la observación sistemática a
gran escala de una zona de acceso público.

El párrafo anterior describe el contenido del artículo pero no lo
transcribe, y no debe citarse como si lo hiciera: su redacción literal y
la de sus apartados relevantes no se han contrastado contra el texto
oficial del Reglamento.

Un sistema de visión computacional operando de forma continua sobre las
zonas comunes de un aeropuerto encaja con naturalidad en esa descripción.

Ahora bien, en este punto concreto la referencia comparada dejó de ser
solo referencia: como se expuso en la Sección~\ref{sec:eipd}, el D.S.
016-2024-JUS establece una obligación propia de Evaluación de Impacto en
Protección de Datos, con supuestos que incluyen expresamente la
videovigilancia masiva. La EIPD peruana \textbf{no} es un estándar
voluntario que el proyecto adopte por prudencia: es derecho aplicable.

El valor que conserva el artículo 35 del RGPD es metodológico. Su
formulación de los supuestos de alto riesgo y de los contenidos de la
evaluación es más detallada y cuenta con abundante doctrina y guías de
autoridades europeas, lo que la hace útil como plantilla de trabajo para
estructurar la EIPD exigida por la norma peruana --- siempre que quede
claro que la obligación se cumple conforme al derecho nacional y que el
RGPD solo aporta el método.

\subsubsection{Reglamento europeo de inteligencia artificial (AI Act)}

El Reglamento (UE) 2024/1689 establece un marco por niveles de riesgo
para los sistemas de inteligencia artificial. Dos elementos son
pertinentes como referencia:

\begin{itemize}[leftmargin=1.5em,itemsep=0.3em]
  \item El régimen prohíbe o somete a condiciones estrictas la
        identificación biométrica remota en tiempo real en espacios de
        acceso público. Un sistema de conteo que no identifica personas
        se sitúa fuera de esa categoría, y esa distinción es
        precisamente el argumento a preservar.
  \item Determinados sistemas quedan clasificados como de alto riesgo
        con obligaciones de gestión de riesgos, gobernanza de datos,
        documentación técnica, registro de eventos, transparencia y
        supervisión humana.
\end{itemize}

Ninguno de los dos puntos anteriores va acompañado de numeral, y la
omisión es deliberada: los artículos concretos que los sustentan y el
calendario de aplicación escalonada del Reglamento no se han contrastado
contra el texto oficial. Ambos se enuncian como descripción del esquema
general de la norma, que es el uso que aquí se les da.

\cajaaviso{Cautela interpretativa}{Sería incorrecto afirmar sin más que
el sistema del proyecto \enquote{no está regulado por el AI Act} o que
\enquote{no es de alto riesgo}. La clasificación depende de la finalidad
declarada y del contexto de uso, no solo de la técnica empleada. Lo que
sí puede sostenerse con base razonable es que un sistema que produce
conteos agregados y no identifica individuos se sitúa en una categoría de
riesgo sustancialmente menor que uno de identificación biométrica.}

\subsubsection{ISO/IEC 27001 e ISO/IEC 42001}

Dos normas técnicas aportan estructura verificable:

\begin{description}[leftmargin=0pt,style=unboxed,itemsep=0.35em]
  \item[ISO/IEC 27001.] Sistema de gestión de seguridad de la
    información. Aporta el marco de controles sobre confidencialidad,
    integridad y disponibilidad que da soporte a las obligaciones de
    seguridad de la Ley 29733 y a la capacidad de detección exigida por
    el régimen de brechas.
  \item[ISO/IEC 42001.] Sistema de gestión de inteligencia artificial.
    Es la norma más reciente y la más específica para este proyecto:
    aporta requisitos de gobernanza sobre el ciclo de vida del modelo,
    incluyendo evaluación de impacto, gestión de datos de entrenamiento y
    supervisión del desempeño en producción.
\end{description}

Ambas se citan aquí por su designación corta, sin año de edición ni
denominación oficial en español, porque esos datos no se han verificado.
Hay además una pregunta previa que conviene hacer al cliente antes de
proponer nada en esta línea: si LAP ya se encuentra certificada en
ISO/IEC 27001. Si lo está, el sistema debería integrarse a su SGSI
existente en lugar de definir un conjunto de controles propio, y la
estrategia de cumplimiento cambia por completo.

\subsection{Privacidad desde el diseño y por defecto: obligación y medidas}
\label{sec:privacidad-diseno}

Conviene corregir de entrada un malentendido frecuente. La privacidad
desde el diseño no es una buena práctica opcional ni un estándar
voluntario importado del régimen europeo: el D.S. 016-2024-JUS la
establece como \textbf{exigencia}, en su doble dimensión de privacidad
\emph{por diseño} --- las garantías se incorporan desde la concepción del
tratamiento, no se añaden después --- y \emph{por defecto} --- la
configuración más protectora es la que debe venir activada de fábrica,
sin requerir que nadie la solicite.

El carácter obligatorio de ambos principios está verificado. Lo que no se
ha contrastado es el artículo que los consagra ni si la norma detalla
criterios o contenidos mínimos para acreditar su cumplimiento; de existir
tales criterios, las medidas que siguen deberían mapearse contra ellos.

La diferencia entre una y otra lectura no es retórica. Si la privacidad
desde el diseño fuese una recomendación, las medidas que siguen serían
propuestas de este frente que el equipo podría descartar por
conveniencia técnica. Al ser una exigencia normativa, son requisitos: su
omisión no es una decisión de producto sino un incumplimiento, y las
decisiones abiertas que se señalan más abajo son decisiones sobre
\emph{cómo} cumplir, no sobre \emph{si} cumplir.

La dimensión «por defecto» tiene además una consecuencia concreta sobre
la configuración del sistema: si una opción admite un ajuste más
protector y otro menos protector --- resolución de captura, frecuencia de
muestreo, granularidad de la salida, plazo de retención ---, el valor de
fábrica debe ser el más protector, y cualquier configuración menos
restrictiva debe ser una desviación deliberada, documentada y
justificada.

Dicho esto, el principio solo es verificable si se traduce en medidas
concretas. Las siguientes se incorporan como requisitos del sistema, con
su fundamento y su forma de verificación.

{\footnotesize
\begin{longtable}{@{}L{1.1cm}L{4.3cm}L{6.6cm}@{}}
\toprule
\textbf{ID} & \textbf{Medida} & \textbf{Fundamento y verificación} \\
\midrule
\endfirsthead
\toprule
\textbf{ID} & \textbf{Medida} & \textbf{Fundamento y verificación} \\
\midrule
\endhead
\bottomrule
\endfoot
M-01 & Procesamiento en el borde: la inferencia se ejecuta en el
terminal y el fotograma no sale del perímetro. &
Minimización y control de transferencia internacional
(Sec.~\ref{sec:extraterritorial}). Verificable por inspección de
arquitectura de red. \\
\addlinespace
M-02 & No persistencia del fotograma: la imagen se procesa en memoria y
se descarta; solo persiste el conteo agregado. &
Es la medida individual con mayor impacto en el perfil de riesgo. Sin
imagen almacenada, la mayor parte del riesgo de brecha desaparece.
Verificable por revisión de código y de esquema de almacenamiento. \\
\addlinespace
M-03 & Salida agregada por zona, sin coordenadas individuales
persistidas. &
Proporcionalidad (Sec.~\ref{sec:directiva}) y límite frente al
seguimiento de trayectorias (Sec.~\ref{sec:ant-p2pnet}). Verificable por
inspección del esquema de la base de datos. \\
\addlinespace
M-04 & Difuminado irreversible de rostros si algún fotograma debe
conservarse para depuración o reentrenamiento. &
Reduce el riesgo de que el material conservado constituya dato
biométrico. Verificable por muestreo del material conservado. \\
\addlinespace
M-05 & Plazo de retención definido y purga automática de los datos
derivados. &
Principio de calidad y de conservación limitada. El plazo concreto no se
fija en este informe: depende de las necesidades analíticas de LAP, que
el equipo aún no conoce. Verificable por configuración de la purga una
vez fijado. \\
\addlinespace
M-06 & Control de acceso por roles y registro de auditoría de todo
acceso. &
Requisito R-01 (Sec.~\ref{sec:brechas}). Verificable por revisión de la
matriz de roles y de los registros. \\
\addlinespace
M-07 & Cifrado en tránsito y en reposo de todo dato derivado. &
Medidas de seguridad exigibles. Verificable por configuración. \\
\addlinespace
M-08 & Carteles informativos en las zonas cubiertas, conforme al modelo
del Anexo~\ref{anx:cartel}. &
Deber de información (Sec.~\ref{sec:directiva}). Verificable por
inspección física. \\
\addlinespace
M-09 & Documentación del linaje de datos: origen, licencia y
transformaciones de cada conjunto usado en entrenamiento. &
Sec.~\ref{sec:licencias}. Verificable por revisión documental. \\
\addlinespace
M-10 & Supervisión humana: el sistema emite alertas, no ejecuta acciones
automáticas sobre personas. &
Evita que una predicción errónea se traduzca directamente en una medida
que afecte a pasajeros. Verificable por diseño del flujo operativo. \\
\end{longtable}}

\cajaaviso{Sobre M-02 y su costo}{M-02 es la medida más protectora y
también la que más restringe al equipo técnico: sin fotogramas
conservados no hay material propio para reentrenar ni para diagnosticar
fallos en producción. La tensión es real y no debe ocultarse. Una salida
intermedia es conservar un subconjunto reducido con difuminado
irreversible (M-04) y plazo corto, bajo un procedimiento documentado.

Cuál de las dos versiones se propone al cliente --- la estricta o la
atenuada --- es una decisión que este informe deja abierta, y conviene
que se tome con los dos costos a la vista: la versión estricta es más
defendible jurídicamente, mientras que la atenuada es la única que
permite mejorar el modelo con datos del propio terminal, que es
justamente lo que exige la brecha de dominio documentada en
\cite{Lin_2025_CVPR}.}

\subsection{Licencias de los conjuntos de datos públicos}
\label{sec:licencias}

El equipo prevé emplear conjuntos públicos de conteo de multitudes. Las
condiciones de uso de estos conjuntos son un punto de cumplimiento
frecuentemente omitido en trabajos académicos y que, en un proyecto con
un cliente real, no puede omitirse.

\cajaaviso{Dónde se hace el seguimiento de esta verificación}{Los
términos de licencia de los conjuntos de conteo de multitudes no siempre
constan en un archivo de licencia estándar; con frecuencia figuran en la
página del grupo de investigación, en un formulario de solicitud o en el
propio artículo. Por esa razón, este informe \textbf{no transcribe} el
texto de ninguna licencia: hacerlo de memoria sería exactamente el tipo
de afirmación no verificable que este documento debe evitar.

La verificación conjunto por conjunto --- ShanghaiTech Partes A y B,
UCF-QNRF y NWPU-Crowd --- está registrada como tarea asignada a
\textbf{Data Engineering} en el archivo
\texttt{CHECKLIST\_\allowbreak CUMPLIMIENTO.md}, sección B. Se ubica ahí
y no en este informe porque requiere acceso a los portales de descarga y
a los formularios de solicitud, que es trabajo del frente de datos y no
del frente ético-legal. Lo que esta sección aporta son los criterios que
esa verificación debe satisfacer y los riesgos que debe descartar.}

Tres riesgos deben quedar señalados desde ahora, con independencia del
resultado de la verificación.

\begin{description}[leftmargin=0pt,style=unboxed,itemsep=0.4em]
  \item[Riesgo de uso comercial.] Es habitual que los conjuntos
    académicos de este dominio se distribuyan para fines exclusivamente
    de investigación no comercial. Si alguno de los tres tuviera esa
    restricción, un modelo entrenado con él no podría desplegarse en un
    entorno operativo de LAP sin analizar previamente el alcance de la
    restricción. Este riesgo no es teórico y debe resolverse
    \textbf{antes} de entrenar, no después.
  \item[Riesgo de transferencia a un tercero en la evaluación.]
    NWPU-Crowd no libera las etiquetas del conjunto de prueba: la
    evaluación se realiza enviando las predicciones a un servidor externo
    administrado por los responsables del conjunto. Esto significa que
    usar la métrica oficial de ese \emph{benchmark} implica
    \textbf{remitir datos a un tercero situado fuera del Perú}. En la
    fase académica lo que se envían son predicciones sobre imágenes del
    propio conjunto público, no material del terminal, de modo que el
    riesgo es acotado; pero el flujo debe quedar documentado, y si en
    algún momento se planteara evaluar de esa forma material propio de
    LAP, se estaría ante una transferencia internacional con todas sus
    consecuencias (Sección~\ref{sec:extraterritorial}).
  \item[Riesgo sobre los titulares de las imágenes.] Estos conjuntos
    contienen fotografías de personas reales que no consintieron a este
    uso específico. Que su empleo sea práctica académica establecida no
    lo convierte en irrelevante desde el punto de vista ético, y así se
    trata en la Sección~\ref{sec:impacto-etico}.
\end{description}

\cajaaviso{Regla propuesta para el servidor de evaluación de NWPU-Crowd}{
Nunca enviar al servidor externo predicciones calculadas sobre imágenes
del Aeropuerto Jorge Chávez, ni ningún artefacto derivado de ellas. El
servidor es una herramienta legítima para comparar contra el estado del
arte sobre datos públicos; deja de serlo en el momento en que el insumo
proviene del terminal. Es una regla fácil de enunciar y fácil de
incumplir por inercia al montar el flujo de evaluación, por lo que
conviene que conste por escrito.}

Si la verificación arroja que alguno de los conjuntos restringe el uso
comercial, se abren dos caminos y el equipo tendrá que elegir uno:
excluir ese conjunto, o emplearlo solo en la fase académica documentando
que el despliegue en LAP requeriría una licencia distinta. La segunda
opción es viable, pero solo si consta por escrito ante el cliente; usarla
sin decirlo generaría la expectativa de que el modelo entregado es
desplegable tal cual, que es precisamente lo que no sería.

\subsection{Auditoría de conformidad de las técnicas del equipo}
\label{sec:auditoria}

Esta subsección cumple la función de auditoría interna: contrasta las
decisiones técnicas del equipo contra los requisitos derivados de las
subsecciones anteriores.

\cajaaviso{Estado de la auditoría}{Al momento de redactar, la
arquitectura definitiva no consta por escrito. La auditoría se formula,
por tanto, de manera condicional sobre los tres enfoques revisados en la
Sección~\ref{sec:antecedentes}, y deberá cerrarse cuando el frente
técnico documente su elección.}

{\footnotesize
\begin{longtable}{@{}L{3.6cm}L{2.0cm}L{6.4cm}@{}}
\toprule
\textbf{Criterio} & \textbf{Estado} & \textbf{Observación} \\
\midrule
\endfirsthead
\toprule
\textbf{Criterio} & \textbf{Estado} & \textbf{Observación} \\
\midrule
\endhead
\bottomrule
\endfoot
EIPD realizada antes de iniciar el tratamiento &
\textbf{No conforme} &
Obligación exigible (Sec.~\ref{sec:eipd}) y aún no cumplida. Es
subsanable de inmediato y con medios propios: la Sec.~\ref{sec:analisis}
ya contiene el análisis sustantivo. \\
\addlinespace
Privacidad por diseño y por defecto &
Parcial &
Las medidas M-01 a M-10 están definidas, pero la dimensión «por defecto»
--- valores de fábrica más protectores --- no está fijada mientras sigan
abiertas las decisiones sobre retención, granularidad y muestreo
(Sec.~\ref{sec:privacidad-diseno}). \\
\addlinespace
No se extraen rasgos biométricos identificantes &
Favorable, condicionado &
Ninguno de los tres métodos revisados construye plantillas faciales. La
conformidad depende de que no se añada un módulo de reidentificación
aguas abajo. \\
\addlinespace
La salida no permite identificar personas &
Favorable con CSRNet; a validar con P2PNet/P2R &
El mapa de densidad es agregado por construcción. Las coordenadas por
punto no identifican por sí solas, pero admiten encadenamiento temporal
(Sec.~\ref{sec:ant-p2pnet}). \\
\addlinespace
Minimización en el proceso de anotación &
Favorable si se adopta P2R &
La eficiencia de etiquetado documentada en \cite{Lin_2025_CVPR} reduce el
volumen de imágenes reales expuestas a anotadores. \\
\addlinespace
Exactitud validada en el dominio de destino &
\textbf{No conforme} &
No existe validación en el terminal. La brecha de dominio documentada en
\cite{Lin_2025_CVPR} impide extrapolar cifras publicadas. Es el hallazgo
más serio de esta auditoría. \\
\addlinespace
Ausencia de sesgo demográfico evaluada &
\textbf{No evaluado} &
Ninguno de los tres artículos reporta desempeño desagregado por
características demográficas. Se desarrolla en
Sec.~\ref{sec:impacto-etico}. \\
\addlinespace
Política de retención definida &
Pendiente &
Depende de M-02 y M-05, ambas con decisión abierta. \\
\addlinespace
Localización del procesamiento &
Pendiente &
Depende de la decisión abierta en Sec.~\ref{sec:extraterritorial}. \\
\addlinespace
Trazabilidad de licencias de datos &
Pendiente &
Verificación asignada a Data Engineering; seguimiento en
\texttt{CHECKLIST\_\allowbreak CUMPLIMIENTO.md}, sección B
(Sec.~\ref{sec:licencias}). \\
\end{longtable}}

Conviene distinguir entre los tres hallazgos no conformes, porque no
tienen la misma naturaleza ni el mismo camino de solución.

La ausencia de EIPD es la más grave en términos formales --- es una
obligación normativa incumplida, no una carencia de evidencia --- pero
también la más fácil de subsanar: no depende de terceros ni de datos que
el equipo no tenga. El material sustantivo ya está redactado en la
Sección~\ref{sec:analisis}; lo que falta es formalizarlo como el
documento que la norma exige. Debería ser lo primero que se cierre.

Los otros dos --- exactitud validada en el dominio de destino y ausencia
de sesgo demográfico evaluada --- no son defectos de implementación sino
ausencias de evidencia, y ambos son subsanables dentro del proyecto: el
primero mediante una campaña de validación en sitio, el segundo mediante
una evaluación de desempeño desagregado. Ambos requieren acceso a datos
del terminal, por lo que no dependen solo del equipo.

Esa dependencia obliga a una gestión concreta: solicitar formalmente a
LAP acceso a material del terminal para la campaña de validación. Es una
petición que el equipo aún no ha cursado al cerrar esta versión del
informe. Si no se obtiene el acceso, los dos hallazgos quedan sin cerrar
y así deben reportarse en las conclusiones, en lugar de presentarse como
resueltos o de omitirse.
